\documentclass[12pt, oneside]{book}

% --- PACKAGES ---
\usepackage{amsmath}       % Math equations
\usepackage{xcolor}        % Colored text
\usepackage{tcolorbox}     % Boxed text/equations
\usepackage{tikz}          % Drawing and image overlays
\usepackage{graphicx}      % Images
\usepackage[colorlinks=true, linkcolor=blue]{hyperref} % Interactive PDF links
\usepackage{parskip}
\usepackage[T1]{fontenc}
\usepackage{amssymb}

\begin{document}

\pagenumbering{roman}
\begin{titlepage}
    \centering
    
    \vspace*{3cm} 
    
    {\Huge \textbf{VQLS Basics}\par}
    
    \vspace*{3cm} 

    \begin{center}
      \hypertarget{maindiagram}{}
      \begin{tikzpicture}
          \node[anchor=south west,inner sep=0] (image) at (0,0) {\includegraphics[width=1.1\textwidth]{VQLSpsn_diagrams/full_diagram_VQLS.png}};
          \begin{scope}[x={(image.south east)}, y={(image.north west)}]
              % PDE Block
            \node at (0.31, 0.80) {\hyperlink{sec:pde}{\phantom{\rule{4cm}{3.10cm}}}};
              % Ansatz Block
              \node at (0.22, 0.44) {\hyperlink{sec:ansatz}{\phantom{\rule{6.35cm}{2.4cm}}}};
              % Cost Function Block
              \node at (0.83, 0.64) {\hyperlink{sec:cost}{\phantom{\rule{4.3cm}{1.8cm}}}};
              % Optimizer Block
              \node at (0.83, 0.2) {\hyperlink{sec:optimizer}{\phantom{\rule{3.6cm}{2.1cm}}}};
          \end{scope}
      \end{tikzpicture}
    \end{center}
    \vfill
    
    {\small Oliver Jackson\par}
    
    \vspace*{0.5cm} 

    {\small \today\par}
\end{titlepage}


\tableofcontents

\cleardoublepage
\pagenumbering{arabic}

\chapter{Introduction}

This document covers variational quantum linear system (VQLS) algorithms where the goal 
is to take a sparse matrix $A$ and a vector $\mathbf{b}$, and prepare a quantum 
state $|x\rangle$ that is proportional to $A^{-1}|b\rangle$.

I hope the reader is familiar with undergraduate linear algebra, systems of 
equations, dirac notation, expectation values, and normalization. Even so, 
I would encourage the inexperienced reader to dive into the document. I have 
provided lots of diagrams, equations, and graphs to help aid the learning process. 

Note if your goal is to see a "quantum advantage" by solving the Poisson 
problem using a VQLS then this project is not for you. In the end, this project 
amounts classically to a small to medium sized matrix multiplication which 
can be handled by your smart phone in a matter of milliseconds to seconds. 
A true advantage may be proposed for quantum systems or large system sizes in which a classical 
computer cannot handle naturally. However, I am interested in solving PDEs 
right now on noisy quantum hardware. So my research is more concerned with 
improving the scalability and runtime issues of VQLS. 

This document was heavily inspired by "Introduction to Variational Quantum 
Algorithms" by Michał Stęchły and "Variational quantum algorithms" by Marco Cerezo et al. 
Michał Stęchły provides an easily digestible non-mathematical introduction to VQA's which 
can give a great base for those coming from non-physics backgrounds. Cerezo's paper 
is one of the original publications on VQA's and provides a comphrehensive overview on the algorithms
and their variations. This document is not a replacement for those two papers but rather a supplement to them 
for those specifically interested in solving PDEs with VQLS. I delve much further into each 
component of the algorithm introducing both the concept and the application. 
\chapter{What are variational quantum linear solvers?}

The VQLS is a hybrid quantum-classical algorithm that uses a parameterized quantum circuit (PQC)
to iteratively prepare a quantum state $|x\rangle$ that is proportional to $A^{-1}|b\rangle$. 
To break that down a little bit, VQLS uses both a quantum processing unit (QPU) and a CPU. It also depends on the 
primitive PQC which for now we can think of as a series of unitary gates of which some depend on 
parameters $\theta$. To find a solution we apply a PQC to a set of quantum states which prepares a guess for $|x\rangle$.
We check how much the guess overlaps with the solution and then update the parameters to make a better guess. You can 
think of VQLS as a quantum analog to classical machine-learning methods like neural networks. That will become 
more apparent as we go through the document. If this was confusing that is ok. In the following sections I break down
the definitions and context of each of these concepts. You can refer back to this definition as you read along. 
\section{Variational quantum algorithms}

VQLS is a type of variational quantum algorithm (VQA) which are an excellent class of algorithms designed to solve problems
on noisy intermediate scale quantum (NISQ) devices. VQAs account for the major limitations of NISQ poor 
coherence times and gate errors by utilizing an optimization or learning based approach. They also have a PQC and optimizer 
but are not limited to solving linear systems of equations. 

VQAs can be used for a wide variety of problems including optimization, machine learning, and most notibly quantum chemistry. 
Of course, whether VQAs can actually provide good solutions to problems on quantum hardware is still an open question. 
The more we research the more we learn about the specific limitations of each application. We are also finding 
that VQAs are not very capable of solving problems on NISQ devices right now but still have potential to be improved. For example,
creating ansatz which are problem specific, optimizers which are robust to noise, and cost functions which are less susceptible
to barren plateaus are all active areas of research. 

All of the accessible quantum hardware to date is noisy and has a limited qubit coherence time and we refer to 
this hardware as NISQ. On IBM's quantum hardware, for example, applying one gate to a qubit introduces an error 
of about 0.1\% to 1\%. This means that if we apply a series of gates to a qubit, the error compounds quickly and the qubit 
will lose its quantum state after only a hundred gates. This limits the depth (number of quantum gates applied) of quantum 
circuits that can be executed on NISQ devices.

So we need an algorithm that can work with shallow circuits. The best known quantum algorithms for solving linear systems 
are the HHL and QSVT algorithms which use the primitives quantum phase estimation (QPE) and quantum singular value transformation
(QSVT) respectively. Both of these primitives require deep circuits and are not suitable for NISQ devices which is why
I have not included them in this project. In interest of seeing if we can solve PDEs on quantum devices today using IBM's 
hardware we will focus on the VQLS algorithm. 

Ok so VQAs are a class of algorithms for NISQ that can solve many problems. But how do they work? How is a VQA constructed? A 
I want to give a little intuition on how to answer these questions before reall diving in. VQA can be described using four steps. 

\begin{enumerate}
  \item Encode the problem 
  \item Choose the ansatz
  \item Build the cost function
  \item Run the optimization loop
\end{enumerate}

The final step which is hidden is to readout the solution. However, this can be a very difficult step and is not always possible. I will 
discuss general readout methods and their limitations in the last chapter. Encoding the problem is the first step and is
when we take our classical problem and encode it into a quantum problem. In a quantum chemistry problem we would 
encode the Hamiltonian of a molecule into a quantum circuit. In the case of a PDE we can encode a matrix $A$ and vector $b$ 
containing the discretized PDE into a quantum circuit. 

For example, we can take any general classical problem represented by the 2x2 matrix 

$$
H = \begin{bmatrix} a & b \\ c & c \\ \end{bmatrix}.
$$

Unlike classical computers, which can store arbitrary matrices as grids of numbers in memory, quantum computers only 
recieve instuctions through a well formulated inner product space. So in order to encode this matrix we have to decompose it into a 
orthonormal basis. The Pauli matrices form an orthonormal basis for 2x2 matrices and consists of the identity matrix and three
Pauli matrices 

$$
I = \begin{bmatrix} 1 & 0 \\ 0 & 1 \\ \end{bmatrix}, \quad 
P_x = \begin{bmatrix} 0 & 1 \\ 1 & 0 \\ \end{bmatrix}, \quad
P_y = \begin{bmatrix} 0 & -i \\ i & 0 \\ \end{bmatrix}, \quad
P_z = \begin{bmatrix} 1 & 0 \\ 0 & -1 \\ \end{bmatrix}. \quad
$$

We could use any orthonormal basis to decompose our matrix into but I have chosen Pauli matrices because they are native to 
IBMs quantum hardware. We can define a new matrix $H$ that is a linear combination of Pauli matrices.

$$
H = c_1X + c_2Y + c_3Z + c_4I 
$$


Where the coefficients $c_k$ are determined by the equation $\frac{1}{2} Tr (A \cdot \sigma_k)$. The ansatz is the specific
architecture you use for your PQC. Mathematically the ansatz is a series of parameterized and unparameterized
unitary gates applied to a quantum state. And it controls how guesses are made for your solution state. In the PQC section
and the ansatz section I discuss exactly how the ansatz prepares a guess solution. 

After choosing our ansatz we need define our cost function. The cost function takes in our encoded matrix as well as the ansatz. 
We define a cost function like in classical machine learning methods to be the discrepancy between our guessed state and our
target state. The final output of the cost function is a real number between 0 and 1. The core goal of the VQA is to minimize this cost 
through optimization. The final step is to choose an optimizer type like gradient descenent. 


\section{Quantum states}
Before jumping into the PQC I want to introduce quantum states and the variational principle to give some motivation to 
the reader on why this algorithm works. Very important. On a classical computer an $n$-bit register lives in a discrete space where each bit can be exactly 
zero or one. In quantum computing, we take the discrete classical states a system can occupy (such as 00,01,10 or 11), make each one an orthogonal axis, 
and assign a complex number to each one. A single quantum state is simply a continuous vector pointing somewhere inside this space.

$$
|\psi\rangle=\sum_{x=0}^{2^n-1} c_x|x\rangle \quad \text{where} \quad|\psi\rangle \in \mathbb{C}^{2^n}
$$

So now our quantum state which is a vector can be formed with any linear combination of possible states and their corresponding 
comlex coefficients. As long as the sum of coefficients squared magnitudes is equal to exactly one. This forms a vector space or more specifically 
an inner product space commonly referred to as the Hilbert space. Therefore, you can mathematically encode $2^n$ classical data points 
into the coefficients of a single $n$-qubit quantum state. Confirm that this makes sense for yourself. 

If a quantum state is a vector pointing somewhere in a continuous complex space, then a quantum logic gate is simply a mathematical 
operation that rotates or transforms that vector. It is a unitary transformation which means it preserves the inner
product (the length) of the state vector. This ensures the total probability of all possible outcomes remains exactly one. 
A transformation is unitary if its conjugate transpose times itself returns the identity matrix 

$$
U^{\dagger} U = I
$$

So a quantum gate can be mathematically represented as a unitary transformation matrix. Let's take the example of an $X$ gate
being applied to one qubit in the zero state. The $X$ gate is represented by the Pauli $x$ matrix and it takes the qubit 
state and flips it. So the transformation is given by 

$$
X |0 \rangle = \begin{bmatrix} 0 & 1 \\ 1 & 0 \\ \end{bmatrix} |0 \rangle = |1 \rangle. 
$$

This is a perfect explanation for what is happening algorithmically, however, it is easy for beginners to confuse actual quantum 
gates with matrices. On real hardware a transformation applied to a qubit by injecting a very specific amount of energy into the 
qubit (such as a photon or electron) which is extremely difficult to do for systems with more than 
one particle. This is why many quantum computing software packages, like qiskit \colorbox{lightgray}{Statevector}, have support for 
qubit simulation. In \colorbox{lightgray}{Statevector} the qubits are manipulated using pauli matrices through 
matrix multiplication. Then by changing our dependence to real hardware like \colorbox{lightgray}{ibm\_boston} those 
same matrices are compiled and translated into machine instructions. 

These instructions are described by the time-dependent schrodinger equation. It is not necessary to understand what it does. 
But notice that there is a dependence on time, and that this equation describes how a quantum state evolves using the 
transformation Hamiltonian $H$ 

$$
i \hbar \frac{d}{d t}|\Psi(t)\rangle=H|\Psi(t)\rangle.
$$

When this equation is solved it has the form

$$
|\Psi(t)\rangle=e^{-i H t}|\Psi(0)\rangle.
$$

This is telling us that any quantum state evolves by taking an initial state and act on it by applying the transformation 
through $e^{-i H t}$ we will get a final state. Now using a similar form we can describe any unitary matrix as an 
imaginary exponential of a Hermitian matrix $H$ 

$$
U = e^{-i H}.
$$

Note that we no longer have dependence on time. Any quantum circuit is built by sequentially applying gates to a qubit
register. Because we know gates are just unitary matrices, a quantum circuit is just the product of these unitary gates 

$$
\text{QC} = \prod_{l=1}^{L} U_l
$$

where $l$ is the number of layers. 

\section{The variational principle}
The variational principle will help us understand why VQAs work. In nature, quantum states are described by the schrodinger 
equation. And the variational principle seeks out only stationary states. A stationary state is an energy eigenstate where 
the energy remains constant. Now take, for example, a single particle in a stationary state. 
Its energy is described by the time-independent form of the schrodinger equation, like so 

$$
\hat{H}|\psi\rangle=E|\psi\rangle.
$$

This is also called the eigenvalue equation in linear algebra. We can find the energy of a specific eigenstate 
by applying the Hamiltonian $H$

$$
\hat{H}|\psi_{\lambda}\rangle=\lambda|\psi_{\lambda}\rangle.
$$

 We might be 
interested in the lowest energy state (the ground state) denoted by 

$$
\hat{H}|\psi_0\rangle = E_0|\psi_0 \rangle.
$$

Once we have the ground state we can determine lots of useful things about our system like its physical or chemical properties.
However, the ground state is not limited to nature and can be reformulated to represent the solution of problems in 
optimization, economics, or PDE solving. We can extract the solution by multiplying the left side with $\langle \psi |$. 
This is the formula for an expectation value of $H$ with respect to $|\psi_0\rangle$. It represents the 
expected value of the ground state energy.   

$$
\langle \psi_0|\hat{H}|\psi_0\rangle = E_0
$$

However, we usually don't know what the specific eigenstates $|\psi_{\lambda}\rangle$ the particle can take on. Because 
the ground state has the absolute lowest energy eigenvalue we know any arbitrary state must have an energy that is 
larger than the ground state. 

$$
\langle \psi | \hat{H} | \psi \rangle \geq E_0
$$

This equation is called the variational principle. It says that the expectation value of a Hamiltonian must be larger than
or equal to the ground state energy. We have a system that can take on a whole space of states, we are looking for a specific eigenstate 
within this space to find its eigenvalue. Now we just need to shuffle through the whole space of potential eigenstates right? 
For large system this can be like finding a needle in a haystack. Instead, we need a tool that will help us quickly explore the 
space of good eigenstates which bring us close to the ground state. And that does not waste any time trying out poor solutions 
that don't bring us close to the ground state. This is where the PQC comes in handy.  


\section{The heart of the VQLS: the parameterized quantum circuit}
A PQC is the same as a quantum circuit except it has parameters 

$$ 
\text{PQC} = \prod_{l=1}^{L} U_l(\theta).
$$

How we define or select these parameters will determine how the PQC explores potential eigenstates. Let's say that our particle 
from earlier only has real eigenstates. We know that a quantum state lives in complex space. So it would be efficient to 
limit that space to only have real values which our solution is also restricted to. Think of this space as a sphere with 
a robot arm attached in the middle. The robot arm can only reach segments of the spheres surface based on what 
gears we give it. Each gear is like a degree of freedom allowing the robot arm to cover more area along the sphere. If we give
the robot arm a gear for an arbitray axis then it could reach all points along that axis. 

\begin{center}
    \includegraphics[width=0.45\textwidth]{VQLSpsn_diagrams/robot_arm.png}
\end{center}

We can give that freedom to our PQC by applying the Pauli matrices from earlier. They will act as our gears. To achieve similar
rotation that the robot arm above does we can use the Pauli y matrix from earlier to formulate a quantum gate given by

$$
R_y(\theta)=e^{-i \frac{\theta}{2} Y}.
$$

Using this form we can construct a simple one qubit PQC 

$$
e^{-i\frac{\theta}{2} Y}|0\rangle = |\psi \rangle.
$$

We can use the other Pauli matrices to add new potential eigenstates to $|\psi\rangle$ giving our circuit more expressivity.  

\chapter{Building a VQLS}

Now that we have the motivation and the tool to solve the problem, we can begin building our VQLS. In the following sections
I break down the four components of a VQLS. There are so many ways to build a VQLS. I have only detailed these sections for 
the application of VQLS for the Poisson problem. This is especially obvious for the first section about encoding the operators. 
There are plenty of books out there purposed for how to discretize PDEs so if you are interested in equations other than the 
Poisson then you will have to apply those texts here. In the rest of the sections they are generalizable to other problems 
you will just have to find the application yourself. 

I have also built this clickable diagram on the cover of the book which outlines the VQLS algorithm in a visual manner. You
can click each of the four main components of the diagram to bring you to the relevant chapter. For example, 
clicking the blue ansatz box with bring you to the ansatz section. Once there you can click the ansatz diagram to 
bring you back to the cover of the document. The diagram corresponds with the material in the following chapters and 
I hope after finishing this document you can use the diagram to see how each of the components fit together. 


\section{Constructing the operators}
\hypertarget{sec:pde}{}
\begin{center}
  \hyperlink{maindiagram}{
    \includegraphics[width=0.65\textwidth]{VQLSpsn_diagrams/pde_diagram_VQLS.png}
  }\par
\end{center}

In my project \colorbox{lightgray}{VQLSpsn} I am solving the 2D Poisson equation. So our goal for this section is to 
construct a Pauli string $\tilde{A}$ which represents the derivative information from our system of equations. I am going to 
first construct a 1D laplacian and decompose into Pauli matrices and then build the 2D laplacian with the 1D laplacian terms. 
This way the number of Pauli strings we end up evaluating is less than if we directly decomposed the 2D laplacian. You will 
see that this means we require less expectation values to complete a training loop. The 1D Poisson laplacian with zero dirichlet
boundary conditions for $N$ qubits is $-2$ along the diagonal and $1$ along the off diagonal terms. For $N=1$ per dimension
we have two points along each axis so the grid spacing is given by $1/2$ and the 1D laplacian is given by 

$$
L_{1D}=\frac{1}{h^2}\cdot \left[\begin{array}{cc}
-2 & 1 \\
1 & -2
\end{array}\right]
=\left[\begin{array}{cc}
-8 & 4 \\
4 & -8
\end{array}\right].
$$

We can decompose this matrix into a linear combination of Pauli matrices. The coefficients are computed using the equation
$\frac{1}{2}Tr[L_{1D} * P]$ where $P$ is a Pauli matrix. We end up getting 

$$
L_{1D}=-8 I+4 X.
$$

The 2D laplacian is given by the equation

$$
A=\left(L_{1D} \otimes I\right)+\left(I \otimes L_{1D}\right).
$$

Then if we multiply out the coefficients and write $A$ as a linear combination of Pauli matrices we get 

$$
\tilde{A}=-16(I I)+4(I X)+4(X I).
$$

This is the final form of our $A$ operator. So you see from the diagram that the final output from this step of the algorithm
is the sum of Pauli operators and their coefficients in $\tilde{A}$ as well as the vector $b$. 


\section{Choosing an ansatz}
\hypertarget{sec:ansatz}{}
\begin{center}
  \hyperlink{maindiagram}{
    \includegraphics[width=0.75\textwidth]{VQLSpsn_diagrams/ansatz_diagram_VQLS.png}
  }\par
\end{center}

If you remember from earlier, an ansatz is just a PQC that also has unparameterized gates. Ansatz prepares a 
quantum state which depends on parameters $\theta$. When we first begin training our VQLS the final state produce results 
that don't align with our solution. However, after many iterations, a properly built VQLS will converge to 
a quantum state which represents an approximation to our solution. This process is guided by the cost function. 
The ansatz takes in initial qubits states $\rho_n$ as well as the parameters from the rotation gates and outputs 
a final quantum state $|x(\theta)\rangle.$ 

$$ 
|x(\theta)\rangle=U(\theta)|0\rangle.
$$

Here $U(\theta)$ is the ansatz being applied to the zero state. The ansatz is defined by 

$$
U_l\left(\theta_l\right)=\prod_m e^{-i \theta_m H_m} W_m.
$$

where $W_m$ is the unparameterized unitary and $H_m$ is the Hermitian operator. In qiskit the ansatz 
\colorbox{lightgray}{real\_amplitudes} has the form 

$$
U_{RA}(\theta)=\prod_m e^{-i \frac{\theta_m}{2} Y_m} \text{CNOT}_m.
$$

where the Hermitian matrix is just the Pauli y matrix and the unparameterized unitary is the CNOT gate. If you 
are unfamiliar with the CNOT gate I would read through qiskit's documentation. When \colorbox{lightgray}{real\_amplitudes} 
is called for two qubits 0 and 1 

$$
U_{RA}(\theta)=\left(\text{CNOT}_{0,1}\right) \cdot\left(e^{-i \frac{\theta_1}{2} Y_1}\right) \cdot\left(e^{-i \frac{\theta_2}{2} Y_2}\right).
$$

Because the solution vector for our problem only has real values I have chosen this to be my ansatz. That way I know my circuit
has enough expressivity but not too much that we are wasting computational resources. If we use Penrose graphical notation 
to visualize this circuit it would have two rotation y gates and one CNOT gate. 



\section{Building the cost function}
\hypertarget{sec:cost}{}
\begin{center}
  \hyperlink{maindiagram}{
    \includegraphics[width=0.75\textwidth]{VQLSpsn_diagrams/cost_diagram_VQLS.png}
  }\par
\end{center}

\section{Choosing an optimizer}
\hypertarget{sec:optimizer}{}
\begin{center}
  \hyperlink{maindiagram}{
    \includegraphics[width=0.75\textwidth]{VQLSpsn_diagrams/optimizer_diagram_VQLS.png}
  }\par
\end{center}

\chapter{The readout problem to be continued}

\end{document}
